% page 84 du fac-simile (image p-083 du scan)
% bloc 160.0 x 233.3 mm ; marge gauche 8.0 mm
\pgc{-12.700mm}{0.000mm}{
\l{\cel{13}76}
\l{}
\l{}
\l{\sou{buduaro}. Mikra chambro di la apartamento di damo, refujeyo ele-}
\l{\cel{3}ganta ube el su lokizas kande elu esas sola. - DEFR.}
\l{}
\l{\sou{bufar}. (\sou{netrans.}) Kreskar volumine ad omna direcioni. - DF.}
\l{}
\l{\sou{bufalo}. (\sou{zool.}) Speco di bovo di la regioni tropikala, aklimat-}
\l{\cel{3}ita en Grekia, Italia, e c., plu granda e plu forta kam la}
\l{\cel{3}bovo ordinara, plu sovaja kam lu, e di qua la muzelo esas}
\l{\cel{3}plu larja kam la lua, la korni plu larja e kurvigita addope.}
\l{\cel{3}- L.\cel{1}\sou{bubalus}. - DEFIS.}
\l{}
\l{\sou{bufeto}. Tablo, plad-etajero ube esas servita dishi; kukaji,}
\l{\cel{3}drinkaji, ye la dispono di la invititi, voyajanti. - DEFR.}
\l{}
\l{\sou{bufono}. La persono qua esforcas ridigar per joki grosiera.}
\l{\cel{3}- DEFRS.}
\l{}
\l{\sou{bufro}.\cel{2}Maso polsterizita (o resortoza) destinita ad amortisar}
\l{\cel{3}la shoki. - DER.}
\l{}
\l{\sou{buglo}. Trumpeto kun klavi. - DEFS.}
\l{}
\l{\sou{bugrano}. Dika telo apretita, quan la taliori uzas kom futero}
\l{\cel{3}interna por rigidigar la infra parto di la pantalono, la}
\l{\cel{3}kolumo, la surfaldaji di vesto. - EFIS.}
\l{}
\l{\sou{buho}. (\sou{zool.}) Nokt-ucelo di la familio \textquotedbl{}strigi\textquotedbl{} (forsan pro la}
\l{\cel{3}plum-penacheti qui ornas lua kapo). - L.\cel{1}\sou{buho maximus}. - L.}
\l{}
\l{\sou{bujio}. I. Kandelo de stearino. - II. Organo uzata por acendar}
\l{\cel{3}la gas-mixuro en la explozo-motori. - DEFS.}
\l{}
\l{\sou{buketo}. Asemblajo de flori koliita e kunligita; fasko analoga}
\l{\cel{3}a la aspekto di flor-buketo. - DEFR.}
\l{}
\l{\sou{buklo}. Sorto di ringo o mi-ringo de metalo tra qua pasas axo}
\l{\cel{3}kun halto-pinti, uzata por fixigar la extremajo di rimeno,}
\l{\cel{3}di gurto, di klapo, e c. - EFIRS.}
\l{}
\l{\sou{budmakero}. Profesionisto di la pariado sur la agri di kaval-}
\l{\cel{3}konkurso-kuri. - EF.}
\l{}
\l{\sou{bukoliko}. Qua apartenas a la genero \textquotedbl{}dramato pastorala\textquotedbl{}. -}
\l{\cel{3}DEFIS.}
\l{}
\l{\sou{bulo}. I. Korpo sfera, generale destinita a rular, a rotacar,}
\l{\cel{3}- II. Globeto quan la aero, la gaso quan kontenas liquido,}
\l{\cel{3}formacas kande lu acensas til la surfaco. - EFS.}
\l{}
\l{\sou{bulbo}. Influro tuberkulala, ovoida, globatra, formacita da la}
\l{\cel{3}asembleso di squami o di folii rudimenta. - EFIS.}
}
